\section{Additional Experiments}
\label{app:additional}

This appendix reports additional measurements that complement the
benchmarks in Section~\ref{sec:results}.
We extend the runtime sweep to stronger exact baselines and to a second GPU,
isolate the contributions of streaming and Tensor-Core tiling through
controlled ablations, and evaluate Flash-SD-KDE on a real-data anomaly
detection benchmark.
All measurements use the 16-D pipeline.
Unless stated otherwise, runtimes are reported in milliseconds, exclude
one-time warmup and JIT compilation costs, and use the 16-D Gaussian-mixture
benchmark of Section~\ref{sec:results} with $n_{\text{test}} = n_{\text{train}}/8$.

\subsection{Extended baseline sweep on the RTX A6000}
\label{app:baseline_sweep}

To place Flash-SD-KDE against a fuller set of exact baselines, we extend the
PyKeOps comparison from a single point to the full $n_{\text{train}}$ sweep
and add a \texttt{torch.compile} variant of the GEMM-based PyTorch baseline.
Table~\ref{tab:rebuttal_baseline_sweep} reports runtimes on the RTX A6000.

\begin{table*}[t]
  \centering
  \small
  \setlength{\tabcolsep}{4pt}
  \begin{tabular}{rrrrrrr}
    \toprule
    $n_{\text{train}}$ & $n_{\text{test}}$ &
    sklearn & Torch & \texttt{torch.compile} & PyKeOps & Flash-SD-KDE \\
    \midrule
    $2{,}048$  & $256$  & $30.7$    & $1.1$   & $6.0$  & $1.9$  & $0.7$ \\
    $4{,}096$  & $512$  & $132.6$   & $2.2$   & $2.9$  & $2.5$  & $0.5$ \\
    $8{,}192$  & $1{,}024$ & $509.7$   & $7.5$   & $5.0$  & $3.6$  & $0.5$ \\
    $16{,}384$ & $2{,}048$ & $1{,}895.7$ & $29.0$  & $11.3$ & $6.7$  & $1.1$ \\
    $32{,}768$ & $4{,}096$ & $7{,}292.7$ & $113.6$ & $34.3$ & $21.7$ & $2.8$ \\
    \bottomrule
  \end{tabular}
  \caption{Average per-call runtime (ms) of exact 16-D SD-KDE implementations
           on the RTX~A6000, lower is better.
           Columns are scikit-learn Gaussian KDE, the GEMM-based PyTorch
           SD-KDE, the same PyTorch implementation under
           \texttt{torch.compile}, the PyKeOps LazyTensor SD-KDE, and our
           Flash-SD-KDE.
           Data are drawn from the 16-D Gaussian-mixture benchmark of
           Section~\ref{sec:results}, with $n_{\text{test}} = n_{\text{train}}/8$.
           One-time warmup and JIT compilation costs are excluded.}
  \label{tab:rebuttal_baseline_sweep}
\end{table*}

Flash-SD-KDE is the fastest exact method at every problem size in the sweep.
At $n_{\text{train}} = 32{,}768$, it is roughly $7.7\times$ faster than PyKeOps,
$12.2\times$ faster than \texttt{torch.compile}, $40.6\times$ faster than the
eager PyTorch baseline, and over three orders of magnitude faster than
scikit-learn.
The PyTorch-based baselines slow down rapidly with $n_{\text{train}}$,
while PyKeOps and Flash-SD-KDE both scale gracefully; among these two,
Flash-SD-KDE retains a consistent margin throughout the sweep.


% \begin{revadded}
\paragraph{Single-point PyKeOps reference.}
The following single-point PyKeOps KDE/SD-KDE comparison was moved from the main text because Table~\ref{tab:rebuttal_baseline_sweep} is the stronger evidence: it compares PyKeOps, \texttt{torch.compile}, eager PyTorch, scikit-learn, and Flash-SD-KDE across the full sweep.  We keep the single-point table here because it directly separates PyKeOps KDE and PyKeOps SD-KDE at the commonly reported $32\text{k}/4\text{k}$ setting.
% \end{revadded}

% [RED WRAP MOVED PYKEOPS SINGLE-POINT TABLE]
% \begin{revadded}
To compare against a specialized KDE implementation, we also benchmarked PyKeOps
LazyTensor operators at $n_{\text{train}}=32\text{k}$, $n_{\text{test}}=4\text{k}$.
PyKeOps currently represents the state of the art for large-scale kernel
reductions, so we report both its Gaussian KDE runtime and its SD-KDE runtime as
a reference point.  Table~\ref{tab:pykeops} shows that Flash-SD-KDE runs
$1.57\times$ faster than the PyKeOps 16-D KDE baseline and $7.99\times$ faster
than the PyKeOps 16-D SD-KDE implementation.

\begin{table}[h]
  \centering
  \small
  \setlength{\tabcolsep}{4pt}
  \begin{tabular}{lcc}
    \toprule
    Method & Runtime (ms) & Rel. to Flash-SD-KDE \\
    \midrule
    16-D Flash-SD-KDE & $2.11$ & $1\times$ \\
    PyKeOps 16-D KDE   & $3.33$ & $1.58\times$ \\
    PyKeOps 16-D SD-KDE& $16.91$ & $8.01\times$ \\
    \bottomrule
  \end{tabular}
  \caption{Runtime comparison at $n_{\text{train}}=32\text{k}$,
           $n_{\text{test}}=4\text{k}$.  PyKeOps rows report Gaussian KDE and
           SD-KDE (LazyTensor) runtimes, while Flash-SD-KDE executes the full SD-KDE
           pipeline.}
  \label{tab:pykeops}
\end{table}
% \end{revadded}
% [END RED WRAP MOVED PYKEOPS SINGLE-POINT TABLE]

\subsection{Generalization across GPU families}
\label{app:hardware}

Section~\ref{sec:results} measures performance on the RTX~A6000.
To check that the gains are not specific to that architecture, we rerun the
same sweep on a Tesla V100-PCIE-16GB.
The V100 has a different SFU/FP32 ratio and a smaller Tensor-Core throughput
than the A6000, so it provides an independent stress test of the
implementation.
Table~\ref{tab:rebuttal_v100} reports the result.

\begin{table*}[t]
  \centering
  \small
  \setlength{\tabcolsep}{4pt}
  \begin{tabular}{rrrrrrr}
    \toprule
    $n_{\text{train}}$ & $n_{\text{test}}$ &
    sklearn & Torch & \texttt{torch.compile} & PyKeOps & Flash-SD-KDE \\
    \midrule
    $2{,}048$  & $256$  & $51.60$   & $0.55$  & $65.20$ & $3.29$  & $0.51$ \\
    $4{,}096$  & $512$  & $270.94$  & $1.78$  & $63.74$ & $4.15$  & $0.62$ \\
    $8{,}192$  & $1{,}024$ & $810.39$  & $6.51$  & $54.25$ & $5.51$  & $1.19$ \\
    $16{,}384$ & $2{,}048$ & $3{,}362.95$ & $25.31$ & $59.88$ & $11.19$ & $3.92$ \\
    \bottomrule
  \end{tabular}
  \caption{Average per-call runtime (ms) on a Tesla V100-PCIE-16GB,
           using the same 16-D Gaussian-mixture benchmark and
           $n_{\text{test}} = n_{\text{train}}/8$ protocol as
           Table~\ref{tab:rebuttal_baseline_sweep}.
           Lower is better.
           One-time warmup and JIT costs are excluded.}
  \label{tab:rebuttal_v100}
\end{table*}

The qualitative ordering is preserved: Flash-SD-KDE remains the fastest
method at every measured size on the V100, with margins comparable to those
seen on the A6000.
The eager PyTorch SD-KDE baseline is competitive on the V100 at small
$n_{\text{train}}$ but degrades rapidly with $n_{\text{train}}$, while
\texttt{torch.compile} pays a large warm overhead that the exclusion of JIT
costs alone does not remove.

\subsection{Effect of query batch size}
\label{app:query_batching}

Many KDE workloads operate at small or moderate query batch sizes, in which
case the score pass dominates and the cost of batching across queries is a
secondary effect.
We isolate this by fixing $n_{\text{train}} = 32{,}768$ and sweeping
$n_{\text{test}}$ across four orders of magnitude.
Table~\ref{tab:rebuttal_query_batch} reports the result.

\begin{table*}[t]
  \centering
  \small
  \setlength{\tabcolsep}{4pt}
  \begin{tabular}{rrrrrr}
    \toprule
    $n_{\text{train}}$ & $n_{\text{test}}$ &
    Torch & \texttt{torch.compile} & PyKeOps & Flash-SD-KDE \\
    \midrule
    $32{,}768$ & $4$      & $100.97$ & $32.44$ & $15.64$ & $1.92$ \\
    $32{,}768$ & $16$     & $101.00$ & $32.41$ & $15.43$ & $1.91$ \\
    $32{,}768$ & $64$     & $103.27$ & $33.44$ & $15.48$ & $1.91$ \\
    $32{,}768$ & $256$    & $104.09$ & $33.68$ & $16.11$ & $1.92$ \\
    $32{,}768$ & $1{,}024$  & $106.22$ & $33.69$ & $15.81$ & $1.97$ \\
    $32{,}768$ & $4{,}096$  & $116.81$ & $35.97$ & $15.85$ & $2.11$ \\
    $32{,}768$ & $16{,}384$ & $158.18$ & $43.14$ & $16.79$ & $2.75$ \\
    \bottomrule
  \end{tabular}
  \caption{Average per-call runtime (ms) on the RTX~A6000 as a function of
           the query batch size $n_{\text{test}}$, with
           $n_{\text{train}} = 32{,}768$ held fixed.
           Lower is better.
           Data are drawn from the 16-D Gaussian-mixture benchmark.}
  \label{tab:rebuttal_query_batch}
\end{table*}

Flash-SD-KDE remains in the $1.91$--$2.75$ ms range across a $4{,}096\times$
sweep of $n_{\text{test}}$, while every baseline shows a noticeable rise at
the largest query batch.
This pattern is consistent with the kernel-level analysis in
Section~\ref{sec:method}: the score pass scales as $O(n_{\text{train}}^2)$
and is independent of $n_{\text{test}}$, so when the score dominates
(\textit{i.e.}\ at large $n_{\text{train}}$ relative to $n_{\text{test}}$),
the end-to-end runtime is essentially flat in $n_{\text{test}}$.
Flash-SD-KDE is therefore particularly attractive in workflows that score a
small number of queries against a large training set, such as anomaly or
out-of-distribution scoring.


% \begin{revadded}
\subsection{Kernel launch-parameter sweep}
\label{app:kernel_tuning}

The full launch-parameter sweep is moved here from the main text.  This preserves the reproducibility detail while keeping Section~\ref{sec:results} focused on the dominant design choices and measured ablations.

Nsight Systems traces show that roughly $95\%$ of the SD-KDE runtime at
$n_{\text{train}}=32\text{k}$, $n_{\text{test}}=4\text{k}$ is spent computing
the empirical score.  We therefore concentrated optimization effort on the
score kernel, sweeping the launch parameters to raise utilization.  Specifically
we varied:
\begin{itemize}
  \item $BLOCK\_M \in \{32, 64, 128, 256\}$,
  \item $BLOCK\_N \in \{32, 64, 128, 256, 512, 1024\}$,
  \item $num\_warps \in \{1, 2, 4, 8\}$,
  \item $num\_stages \in \{1, 2, 4\}$,
\end{itemize}
and selected the combination that minimized runtime.  For this workload we
found that $BLOCK\_M=64$, $BLOCK\_N=1024$, $num\_warps=2$, and $num\_stages=2$
gave the best overall performance, increasing measured FLOP utilization by
more than $2\times$ relative to smaller-tile settings.
% \end{revadded}


\subsection{Streaming versus full materialization}
\label{app:streaming_ablation}

The streaming design described in Section~\ref{sec:method} avoids ever
materializing the full pairwise kernel matrices.
To quantify how much of the speedup is attributable to streaming alone, we
compare the streamed Flash-SD-KDE kernel against an explicit
``full-materialization'' baseline that allocates the
$n_{\text{train}} \times n_{\text{train}}$ training kernel matrix and the
$n_{\text{test}} \times n_{\text{train}}$ query kernel matrix.
Both implementations otherwise share the same numerical formulation.
Table~\ref{tab:rebuttal_materialization} reports runtime and peak GPU
memory above the resident input-tensor footprint at the largest size we can
materialize on a single A6000, $n_{\text{train}} = 65{,}536$,
$n_{\text{test}} = 8{,}192$.

\begin{table*}[t]
  \centering
  \small
  \setlength{\tabcolsep}{4pt}
  \begin{tabular}{lrr}
    \toprule
    Method & Runtime (ms) & Peak alloc (MB) \\
    \midrule
    Streamed Flash + Tensor Cores      & $10.99$  & $16.25$       \\
    Streamed Flash + No Tensor Cores   & $49.52$  & $16.25$       \\
    Full materialization + Tensor Cores& $611.63$ & $16{,}400.50$ \\
    Full materialization + No Tensor Cores & $617.49$ & $16{,}400.50$ \\
    \bottomrule
  \end{tabular}
  \caption{Runtime and peak extra GPU memory at $n_{\text{train}} = 65{,}536$,
           $n_{\text{test}} = 8{,}192$ on the RTX~A6000.
           ``Streamed Flash'' is the production Flash-SD-KDE kernel;
           ``Full materialization'' explicitly forms the
           $n_{\text{train}} \times n_{\text{train}}$ and
           $n_{\text{test}} \times n_{\text{train}}$ kernel matrices in
           global memory.
           Peak alloc is reported above the input-tensor footprint and
           excludes shared workspace.}
  \label{tab:rebuttal_materialization}
\end{table*}

Streaming reduces peak extra GPU memory by roughly $1{,}000\times$ at this
size, from $\sim\!16$ GB to $\sim\!16$ MB, and gives a $55.7\times$ runtime
speedup over the materialized Tensor-Core variant.
The materialized baseline cannot scale meaningfully beyond this size on the
A6000, while streaming Flash-SD-KDE continues to operate well under typical
GPU memory budgets.
Streaming therefore accounts for the bulk of the memory savings reported in
Section~\ref{sec:results}, separately from the Tensor-Core speedup studied in
Section~\ref{app:tc_ablation}.

\subsection{Tensor-Core ablation within the streamed kernel}
\label{app:tc_ablation}

Within the streamed Flash-SD-KDE kernel we additionally compare a Tensor-Core
GEMM path against an otherwise identical FP32 path that disables Tensor-Core
tiling.
This isolates the Tensor-Core contribution from streaming and from the
launch-parameter tuning of Section~\ref{sec:method}.
Table~\ref{tab:rebuttal_tc_ablation} sweeps $n_{\text{train}}$ from $2{,}048$
to $65{,}536$ with $n_{\text{test}} = n_{\text{train}}/8$.

\begin{table}[h]
  \centering
  \small
  \setlength{\tabcolsep}{4pt}
  \begin{tabular}{rrrrr}
    \toprule
    $n_{\text{train}}$ & $n_{\text{test}}$ &
    TC (ms) & no-TC (ms) & no-TC / TC \\
    \midrule
    $2{,}048$  & $256$    & $0.30$ & $0.31$  & $1.04\times$ \\
    $4{,}096$  & $512$    & $0.28$ & $0.35$  & $1.24\times$ \\
    $8{,}192$  & $1{,}024$  & $0.29$ & $0.87$  & $2.99\times$ \\
    $16{,}384$ & $2{,}048$  & $0.66$ & $2.83$  & $4.27\times$ \\
    $32{,}768$ & $4{,}096$  & $2.12$ & $10.48$ & $4.94\times$ \\
    $65{,}536$ & $8{,}192$  & $9.30$ & $44.59$ & $4.80\times$ \\
    \bottomrule
  \end{tabular}
  \caption{Tensor-Core ablation within the streamed Flash-SD-KDE kernel on
           the RTX~A6000.
           ``TC'' is the production kernel; ``no-TC'' is an otherwise
           identical kernel with Tensor-Core tiling disabled.
           The right-most column reports
           $\text{runtime}_{\text{no-TC}}/\text{runtime}_{\text{TC}}$, so
           values above $1\times$ mean the Tensor-Core path is faster.
           Lower runtime is better.}
  \label{tab:rebuttal_tc_ablation}
\end{table}

The Tensor-Core advantage grows from a marginal $1.04\times$ at the smallest
size to $4.94\times$ at $n_{\text{train}} = 32{,}768$, and stabilizes near
$5\times$ at the largest sizes we can run.
This is consistent with arithmetic-intensity reasoning: at small
$n_{\text{train}}$ the GEMM tiles are too small to saturate Tensor Cores,
while at moderate-to-large $n_{\text{train}}$ the GEMM portion dominates and
the kernel approaches the mixed Tensor-Core/FP32 roofline analysed in
Section~\ref{sec:method}.
Together with Section~\ref{app:streaming_ablation}, this gives a clean
attribution: streaming is responsible for the memory-footprint reduction at
all sizes, while Tensor-Core tiling drives the runtime gain at scale.

We do not currently fuse the score pass with the KDE pass into a single
kernel.
Nsight Systems traces show that the score pass accounts for roughly
$90\%$ of the end-to-end runtime, so fusing the two kernels yields only a
small incremental gain at the cost of substantial additional code complexity.
We retain the simpler two-kernel structure.

\subsection{Anomaly detection on the ACE benchmark}
\label{app:ace}

The benchmarks above measure runtime and oracle accuracy.
To check that Flash-SD-KDE is competitive on a real, non-synthetic task, we
evaluate it on the anomaly-detection benchmark of \citet{luo2018ace}.
We follow the same protocol: rank all points by an outlier score, threshold
at the top-$K$ to recover a candidate outlier set, and report the number of
true outliers correctly recovered, the total number reported, the number
missed, and the wall-clock execution time.
Flash-SD-KDE produces an outlier score by negating its density estimate, so
that low-density points rank as outliers.

We compare against the twelve baselines reported in the ACE paper on two
of the three datasets evaluated there: the Statlog Shuttle dataset
($n=34{,}987$, $d=9$, $879$ outliers) and the ALOI object-image dataset
($n=50{,}000$, $d=27$, $1{,}508$ outliers).
The non-Flash-SD-KDE rows in
Tables~\ref{tab:rebuttal_ace_shuttle}--\ref{tab:rebuttal_ace_aloi} are taken
directly from the ACE paper; the Flash-SD-KDE rows are our measurements on
the RTX~A6000 using the same datasets and the same top-$K$ protocol.

\begin{table*}[t]
  \centering
  \small
  \setlength{\tabcolsep}{4pt}
  \begin{tabular}{lrrrrr}
    \toprule
    Method        & Reported & Correct & Missed & Time (s) & vs ACE \\
    \midrule
    ACE           & $6{,}763$  & $273$ & $606$ & $0.81$    & $1\times$    \\
    LOF           & $4{,}356$  & $381$ & $498$ & $14.12$   & $17.4\times$ \\
    kNN           & $4{,}897$  & $493$ & $386$ & $12.35$   & $15.2\times$ \\
    kNNW          & $5{,}264$  & $610$ & $269$ & $13.54$   & $16.7\times$ \\
    LoOP          & $6{,}145$  & $201$ & $678$ & $14.51$   & $17.9\times$ \\
    LDOF          & $6{,}433$  & $330$ & $549$ & $16.42$   & $20.3\times$ \\
    ODIN          & $9{,}775$  & $375$ & $504$ & $12.21$   & $15.1\times$ \\
    KDEOS         & $12{,}630$ & $314$ & $565$ & $11.73$   & $14.5\times$ \\
    COF           & $9{,}133$  & $280$ & $599$ & $13.45$   & $16.6\times$ \\
    LDF           & $9{,}809$  & $375$ & $504$ & $19.93$   & $24.6\times$ \\
    INFLO         & $4{,}488$  & $183$ & $696$ & $14.03$   & $17.3\times$ \\
    FastVOA       & $8{,}532$  & $271$ & $608$ & $235.10$  & $290.2\times$ \\
    Flash-SD-KDE  & $5{,}130$  & $807$ & $72$  & $0.60$    & $0.74\times$ \\
    \bottomrule
  \end{tabular}
  \caption{Anomaly recovery on Statlog Shuttle ($n=34{,}987$, $d=9$,
           $879$ true outliers).
           ``Reported'' is the number of points returned at the chosen
           threshold, ``Correct'' the number of those that are true outliers,
           and ``Missed'' the number of true outliers not recovered.
           Time is total execution time in seconds; ``vs ACE'' is the
           runtime ratio relative to the ACE row.
           Non-Flash-SD-KDE rows are reproduced from the ACE paper;
           Flash-SD-KDE is measured on the RTX~A6000.}
  \label{tab:rebuttal_ace_shuttle}
\end{table*}

\begin{table*}[t]
  \centering
  \small
  \setlength{\tabcolsep}{4pt}
  \begin{tabular}{lrrrrr}
    \toprule
    Method        & Reported & Correct & Missed  & Time (s) & vs ACE \\
    \midrule
    ACE           & $7{,}216$  & $340$ & $1{,}168$ & $1.26$   & $1\times$    \\
    LOF           & $4{,}476$  & $519$ & $989$     & $72.31$  & $57.4\times$ \\
    kNN           & $5{,}428$  & $447$ & $1{,}061$ & $63.27$  & $50.2\times$ \\
    kNNW          & $5{,}558$  & $329$ & $1{,}508$ & $89.96$  & $71.4\times$ \\
    LoOP          & $5{,}121$  & $253$ & $1{,}179$ & $59.97$  & $47.6\times$ \\
    LDOF          & $7{,}501$  & $470$ & $1{,}038$ & $60.39$  & $47.9\times$ \\
    ODIN          & $10{,}110$ & $162$ & $1{,}346$ & $72.69$  & $57.6\times$ \\
    KDEOS         & $9{,}515$  & $404$ & $1{,}104$ & $55.89$  & $44.4\times$ \\
    COF           & $8{,}746$  & $284$ & $1{,}224$ & $81.74$  & $64.9\times$ \\
    LDF           & $9{,}133$  & $301$ & $1{,}207$ & $60.51$  & $48.0\times$ \\
    INFLO         & $10{,}328$ & $420$ & $1{,}088$ & $72.13$  & $57.2\times$ \\
    FastVOA       & $8{,}931$  & $319$ & $1{,}189$ & $291.10$ & $231.0\times$ \\
    Flash-SD-KDE  & $7{,}250$  & $437$ & $1{,}071$ & $0.09$   & $0.07\times$ \\
    \bottomrule
  \end{tabular}
  \caption{Anomaly recovery on the ALOI object-image dataset
           ($n=50{,}000$, $d=27$, $1{,}508$ true outliers).
           Columns and source attribution match
           Table~\ref{tab:rebuttal_ace_shuttle}.}
  \label{tab:rebuttal_ace_aloi}
\end{table*}

On Shuttle, Flash-SD-KDE recovers $807$ of the $879$ true outliers at the
chosen threshold, compared with $610$ for the strongest non-Flash baseline
(kNNW) and $273$ for ACE itself, while running in $0.60$ s versus the next
fastest method, ACE at $0.81$ s.
On ALOI, Flash-SD-KDE is competitive on recovery quality (within $20\%$ of
the best non-Flash baseline) and is more than an order of magnitude faster
than ACE, and roughly three orders of magnitude faster than the
distance/density-based baselines.
The two datasets together cover relatively low ($d=9$) and moderate ($d=27$)
dimensionalities, both of which are within the regime where the Tensor-Core
formulation of Section~\ref{sec:method} is effective.

We additionally evaluate Flash-SD-KDE at large scale on a public point-anomaly
release of the KDD-Cup~99 HTTP dataset ($n=620{,}098$, $d=29$, $1{,}052$
labelled anomalies), which is comparable in scale to the
$596{,}853$-point KDD HTTP subset used in the ACE paper.
On this release Flash-SD-KDE reaches a ROC-AUC of $0.92$ in $7.27$ s of
end-to-end execution time on the RTX~A6000.
For reference, the ACE paper reports $23.33$ s on its own KDD HTTP subset.
We do not claim a like-for-like comparison here because the public release
does not match the exact subset used in the ACE paper, but the result
indicates that Flash-SD-KDE remains practical at the $\sim\!10^6$-point
scale on commodity hardware.
